\clearpage\thispagestyle{empty}
\pdfbookmark[0]{Quellen zu Pallas (Briefe vom 1. Februar und 24. April 1815)}{pallascorr1815-sources}
\begin{center}\textbf{Quellen zu Pallas (Briefe vom 1. Februar und 24. April 1815)}\end{center}
\noindent Diese KI-Übertragung ins heutige Deutsch folgt den beiden Briefen von Carl Friedrich Gauß an Johann Elert Bode unter Pallas in den \href{https://gdz.sub.uni-goettingen.de/id/PPN236007467}{\textit{Werken}, Band~VI} (Göttingen, 1874), S.~383--385. Bode druckte die Auszüge unter einer gemeinsamen Überschrift zu Pallas, zur Wintersonnenwende von 1814, zum Kometen von 1815 und zur barometrischen Höhenmessung im \textit{Astronomischen Jahrbuch für 1818} (Berlin, 1815), S.~\href{https://api.digitale-sammlungen.de/iiif/image/v2/bsb10538347_00175/full/1600,/0/default.jpg}{167}--\href{https://api.digitale-sammlungen.de/iiif/image/v2/bsb10538347_00181/full/1600,/0/default.jpg}{173}. Die Links führen zum Exemplar der Bayerischen Staatsbibliothek.
\par\smallskip\noindent Die Fußnote zu Enckes Ephemeride auf S.~\href{https://api.digitale-sammlungen.de/iiif/image/v2/bsb10538347_00176/full/1600,/0/default.jpg}{168} lautet „Sie steht schon oben.“ Der Abschnitt zur barometrischen Höhenmessung mit drei Tafeln und einem Rechenbeispiel reicht im Jahrbuch von S.~\href{https://api.digitale-sammlungen.de/iiif/image/v2/bsb10538347_00177/full/1600,/0/default.jpg}{169} bis zum oberen Teil von S.~\href{https://api.digitale-sammlungen.de/iiif/image/v2/bsb10538347_00180/full/1600,/0/default.jpg}{172}. In der Werkausgabe fehlen beide in Band~VI.
\par\smallskip\noindent Das Jahrbuch gibt für die mittlere Zeit am 25.~März 9~46~58 und für die Rektaszension am 30.~März 59~13~3 an. In den Briefen der Werkausgabe stehen 9~46~38 und 59~14~3. Für die Länge des Perihels nennt das Jahrbuch $146^\circ\,7'\,2''$, die Briefe der Werkausgabe nennen $146^\circ\,7\,2$. Im gesonderten Bericht vom 8.~April steht in der Werkausgabe $246^\circ\,7'\,2''$. KI-Übertragung und KI-Satz.
\clearpage
